\chapter{Determining the Chemical Potential from the Available States}
\label{ch:determining-chemical-potential}

An energy band specifies the available one-particle states and their energies.
The partition function gives their occupations for a trial chemical potential.
Particle-number conservation or charge neutrality selects the physical value of
$\mu$.

This chapter develops the general constrained problem before specializing to
intrinsic and extrinsic semiconductors.  The central equation is either a fixed
particle-number condition,
\begin{equation}
  N(\mu,T)=N_{\mathrm{target}},
  \label{eq:general-fixed-particle-constraint}
\end{equation}
or a fixed-charge condition,
\begin{equation}
  Q(\mu,T)=Q_{\mathrm{target}}.
  \label{eq:general-fixed-charge-constraint}
\end{equation}
The chemical potential is the value that satisfies the declared constraint.

\section{Thermodynamic inputs beyond the band structure}

Consider two samples made from the same semiconductor.  If their crystal
structures and temperatures are equal, they can have essentially the same host
bands $E_n(\bm k)$ and density of states $g(E)$.  If one sample is donor doped
and the other acceptor doped, their equilibrium chemical potentials are
different.  Their total-charge constraints and dopant charge states select the
different values of $\mu$.

Even in an undoped material, temperature matters.  The same fixed electron
number is distributed more broadly in energy as $T$ increases, and $\mu(T)$ can
move to preserve that number.  A chemical-potential calculation therefore specifies the temperature, the
conserved particle number or charge, the included states and degeneracies, and a
common energy reference for bands, localized levels, and $\mu$.

\begin{boundarybox}
Band energies determine the possible occupations.  Temperature determines how
sharply those occupations change with energy.  Particle number or charge
neutrality determines which chemical potential is realized.
\end{boundarybox}

\section{The finite discrete-state problem}

Begin with discrete one-particle energies $E_i$ and multiplicities $g_i$.  At
thermal equilibrium, the expected electron number is
\begin{equation}
  \boxed{
  N(\mu,T)
  =\sum_i g_i
   \frac{1}{1+\exp[(E_i-\mu)/(\kb T)]}.
  }
  \label{eq:discrete-electron-number-constraint}
\end{equation}
The multiplicity $g_i$ counts how many distinct one-particle states have energy
$E_i$; the Fermi--Dirac factor gives the occupied fraction of those states.
The equation
\begin{equation}
  N(\mu,T)=N_{\mathrm{target}}
  \label{eq:discrete-chemical-potential-equation}
\end{equation}
then determines $\mu$.

The left side is monotonic.  Differentiating one occupation gives
\begin{equation}
\begin{aligned}
  \frac{\partial f_i}{\partial\mu}
  &=\beta
    \frac{e^{\beta(E_i-\mu)}}
         {[1+e^{\beta(E_i-\mu)}]^2}\\
  &=\beta f_i(1-f_i)>0.
\end{aligned}
  \label{eq:fermi-occupation-mu-derivative}
\end{equation}
Therefore
\begin{equation}
  \frac{\partial N}{\partial\mu}
  =\beta\sum_i g_i f_i(1-f_i)>0
  \label{eq:discrete-particle-number-monotonicity}
\end{equation}
whenever at least one state is partially occupied.  Raising $\mu$ always raises
the expected particle number.  This monotonicity normally makes the solution
unique.

\subsection{Zero-temperature limit}

As $T\rightarrow0$,
\begin{equation}
  f(E,\mu,0)
  =\begin{cases}
      1, & E<\mu,\\
      0, & E>\mu.
    \end{cases}
  \label{eq:zero-temperature-step-occupation}
\end{equation}
States are filled from lowest energy upward.  For a partially filled level,
$\mu$ coincides with that level in the limiting description.  If an integer
number of complete levels is filled and a gap follows, any $\mu$ inside that
gap produces the same zero-temperature occupations.  The finite-temperature
limit or a thermodynamic convention selects a particular limiting value.

A semiconductor chemical potential can lie inside the band gap between the
allowed one-particle states.  It remains the parameter controlling the relative
statistical cost of adding or removing an electron.

\subsection{A symmetric two-level example}

Let a spinless particle occupy one of two levels,
\begin{equation}
  E_-=E_0-\frac{\Delta}{2},
  \qquad
  E_+=E_0+\frac{\Delta}{2},
  \label{eq:symmetric-two-level-energies}
\end{equation}
and require one expected particle:
\begin{equation}
  f(E_-,\mu,T)+f(E_+,\mu,T)=1.
  \label{eq:symmetric-two-level-number-constraint}
\end{equation}
Using
\begin{equation}
  f(E_0-\epsilon,E_0,T)
  =1-f(E_0+\epsilon,E_0,T),
  \label{eq:fermi-symmetry-about-mu}
\end{equation}
the constraint gives
\begin{equation}
  \boxed{\mu=E_0}
  \label{eq:symmetric-two-level-chemical-potential}
\end{equation}
for every $T$.  Spectral and particle--hole symmetry place the chemical
potential at the arithmetic mean in this example.

\section{The continuum limit for energy bands}

For a macroscopic solid, the discrete sum becomes a density-of-states integral.
If all relevant electronic bands are included,
\begin{equation}
  \boxed{
  \frac{N(\mu,T)}{V}
  =\sum_n\int g_n(E)f(E,\mu,T)\,\dd E.
  }
  \label{eq:continuum-total-electron-number}
\end{equation}
The prescribed total electron density determines $\mu$ through
\begin{equation}
  \sum_n\int g_n(E)f(E,\mu,T)\,\dd E
  =\frac{N_{\mathrm{target}}}{V}.
  \label{eq:continuum-electron-number-closure}
\end{equation}

This expression is exact for the declared independent-particle band model, but
it can be inconvenient for a semiconductor.  The total includes the enormous
population of occupied valence states and perhaps deeper core states, while the
quantities of interest are often the much smaller deviations from a neutral
reference.  Subtracting two large numbers to find a small carrier imbalance is
also numerically fragile.

The electron--hole description reorganizes the same counting around a neutral
reference.  Conduction electrons are occupied states above $E_C$:
\begin{equation}
  n(\mu,T)
  =\int_{E_C}^{\infty}g_c(E)f(E,\mu,T)\,\dd E.
  \label{eq:chemical-potential-conduction-electrons}
\end{equation}
Valence holes are missing electrons below $E_V$:
\begin{equation}
  p(\mu,T)
  =\int_{-\infty}^{E_V}g_v(E)[1-f(E,\mu,T)]\,\dd E.
  \label{eq:chemical-potential-valence-holes}
\end{equation}
The fixed positive charge of the host lattice is already incorporated in the
neutral reference.  The remaining constraint can then be written directly in
terms of the mobile carrier and dopant deviations.

\section{Charge neutrality as the semiconductor constraint}

For a uniform semiconductor containing singly charged donors and acceptors,
\begin{equation}
  \rho(\mu,T)
  =q\left[p+N_D^+-n-N_A^-\right].
  \label{eq:chemical-potential-charge-density}
\end{equation}
Bulk neutrality requires
\begin{equation}
  \boxed{
  n(\mu,T)+N_A^-(\mu,T)
  =p(\mu,T)+N_D^+(\mu,T).
  }
  \label{eq:chemical-potential-neutrality-closure}
\end{equation}
This one equation combines four statistical populations that depend on the
same chemical potential.

Define the neutrality residual
\begin{equation}
  R(\mu)
  =n(\mu,T)+N_A^-(\mu,T)
   -p(\mu,T)-N_D^+(\mu,T).
  \label{eq:chemical-potential-neutrality-residual}
\end{equation}
The equilibrium solution satisfies
\begin{equation}
  R(\mu)=0.
  \label{eq:chemical-potential-root-equation}
\end{equation}
Increasing $\mu$ raises the conduction-electron population and the negatively
charged acceptor population.  It lowers the valence-hole population and the
positively charged donor population.  Every change raises $R(\mu)$ under the
independent-state model.
Consequently the root is monotonic and normally unique.

\section{Intrinsic material as the first closure problem}

An ideal intrinsic semiconductor contains no charged dopants:
\begin{equation}
  N_D^+=N_A^-=0.
  \label{eq:no-ionized-dopants-intrinsic}
\end{equation}
Neutrality reduces to
\begin{equation}
  n(\mui,T)=p(\mui,T).
  \label{eq:chemical-potential-intrinsic-neutrality}
\end{equation}
The intrinsic chemical potential $\mui$ is the value at which the statistically
expected conduction-electron and valence-hole populations are equal.  Equal
conduction and valence state counts place it at midgap; unequal state counts
displace it toward the band with fewer thermally accessible states.

Under nondegenerate parabolic-band conditions,
\begin{align}
  n&=N_Ce^{-(E_C-\mu)/(\kb T)},
  \label{eq:chemical-potential-boltzmann-electrons}\\
  p&=N_Ve^{-(\mu-E_V)/(\kb T)}.
  \label{eq:chemical-potential-boltzmann-holes}
\end{align}
Equating them gives
\begin{equation}
  \mui
  =\frac{E_C+E_V}{2}
   +\frac{\kb T}{2}\ln\left(\frac{N_V}{N_C}\right).
  \label{eq:chemical-potential-intrinsic-result-preview}
\end{equation}
The following intrinsic-semiconductor chapter derives and interprets each term.

\section{Energy references must be shared}

Equilibrium occupations depend on energy differences such as $E-\mu$,
$E_C-\mu$, and $E_D-\mu$.  If every electronic energy is shifted by one
constant $C$,
\begin{equation}
  E' = E+C,
  \qquad
  \mu'=\mu+C,
  \label{eq:chemical-potential-rigid-energy-shift}
\end{equation}
then
\begin{equation}
  E'-\mu'=E-\mu
  \label{eq:chemical-potential-energy-difference-invariance}
\end{equation}
and all populations remain unchanged.

Conduction bands, valence bands, localized levels, and $\mu$ require one aligned
energy reference.  Explicit offsets connect conventions based on the vacuum
level, valence-band maximum, average electrostatic potential, or a numerical
energy zero.

\section{Choosing the statistical and band model}

The physical regime determines the band and occupation models used in the
closing equation.

\subsection{Nondegenerate or Fermi--Dirac statistics}

If $\mu$ lies several $\kb T$ from a band edge, the Boltzmann approximation can
be checked against the full Fermi integral.  If $\mu$ approaches or enters a
band, the full Fermi--Dirac occupation is required.  Using a Boltzmann formula
to calculate a chemical potential and then finding that the result lies inside
the band is a failed self-consistency check.

\subsection{Parabolic or full bands}

Near a band edge, effective masses and analytic densities of states may be
adequate.  If thermal occupation samples nonparabolic regions, multiple remote
valleys, or narrow bands, the carrier integrals must use a more complete
$g(E)$ calculated from the bands.

\subsection{Complete or incomplete dopant ionization}

The replacements $N_D^+=N_D$ and $N_A^-=N_A$ are regime-dependent
approximations.  At low temperature or for deeper levels, localized-state
partition functions must be included in the same neutrality equation.

\subsection{Fixed bands or charge-dependent bands}

At high carrier or impurity concentration, screening, exchange, band-gap
narrowing, and impurity-band formation can make the bands themselves depend on
the solution.  The problem is then self-consistent at two levels: occupations
determine charge, and charge modifies the available states.

\section{Spatially varying systems}

A single bulk chemical potential can describe thermal equilibrium, but the
local band edges can vary with electrostatic potential:
\begin{equation}
  E_C(\bm r)=E_C^{(0)}-q\phi(\bm r),
  \qquad
  E_V(\bm r)=E_V^{(0)}-q\phi(\bm r).
  \label{eq:local-band-edges-electrostatic-potential}
\end{equation}
The local carrier populations then depend on $\mu-E_C(\bm r)$ and
$E_V(\bm r)-\mu$.  Their charge density enters Poisson's equation,
\begin{equation}
  -\nabla\cdot[\varepsilon(\bm r)\nabla\phi(\bm r)]
  =\rho[\phi(\bm r),\mu,T].
  \label{eq:chemical-potential-poisson-closure}
\end{equation}
The equilibrium calculation must solve the occupation and electrostatic
problems together.  Imposing local neutrality would remove the space charge
that creates depletion regions and built-in fields.

Outside equilibrium, electron and hole quasi-Fermi levels describe carrier
populations driven by generation, recombination, injection, or external fields.
A single chemical potential describes thermal equilibrium.

\section{A reliable mathematical solution procedure}

For a uniform equilibrium problem, all populations are expressed as functions
of one trial $\mu$ and combined into the residual $R(\mu)$.  The calculation
uses one energy reference, a specified temperature and dimensionality, a
normalized DOS with declared degeneracies, and a complete account of the
conserved particle number or charge.  The selected occupation laws state their
physical regime.  Opposite residual signs at two trial chemical potentials
form a bracket containing the root.

The final value is substituted into the original number or charge equation.
Energy-shift invariance checks the alignment convention, while refinement of
the band integrals checks numerical convergence.  The returned $\mu$ must lie
inside the physical regime of every approximation used to construct the
populations.

The bracket is physically interpretable.  Far below all relevant electronic
states, the system lacks electrons and is positively charged relative to the
target.  Far above them, the states fill and the system becomes negatively
charged.  The equilibrium chemical potential lies between these limiting charge
states.

\section{Physical closure}

The one-electron spectrum and its spectral multiplicities define the normalized
density of states.  The partition function defines the many-electron occupation
of those states.  Valence vacancies define holes relative to the neutral ionic
background, while localized dopant partition sums define additional charge
states.  A specified temperature and a conserved particle number or charge then
produce one closing equation for $\mu$ on the common energy scale.

\section{Chapter summary}

The chemical potential follows from a band spectrum together with thermodynamic
information.  Bands and their multiplicities determine the available states;
the partition function determines their expected occupation
for a trial $\mu$; and particle-number conservation or charge neutrality
selects the physical value.  In semiconductors, electron--hole notation recasts
the total-electron problem relative to a neutral filled-valence reference and
avoids subtracting enormous nearly equal charges.  Exact bands, parabolic
reductions, Fermi integrals, Boltzmann limits, and dopant-ionization
approximations are alternative models for the same constrained calculation.
The following chapters apply this framework first to intrinsic material and
then to extrinsic material.
